\documentclass[12pt]{evaluasep}
\grado{2$^\circ$ de Secundaria}
\cicloescolar{2025-2026}
\materia{Olimpiada de Conocimientos 2026}
\unidad{1}
\title{Propuesta para Examen}
\author{Profesores de Secundaria}
\begin{document}
% \begin{multicols}{2}
	\tableofcontents
% \end{multicols}
\vfill

\newpage
\begin{questions}
	\section{Lenjuages}
                \subsection{Lengua Materna}

\question ¿Cuáles son considerados los tres tipos fundamentales de violencia y acoso sexual?
    \begin{choices}
        \choice Físico, cibernético y verbal de carácter informal.
        \CorrectChoice Físico, psicológico y sexual.
        \choice Psicológico, bullying escolar y físico directo.
        \choice Sexual, psicológico y laboral o corporativo.
    \end{choices}
    \vspace{0.3cm}

    \question Indica cuáles son las tres etapas fundamentales en la estructura de un debate formal:
    \begin{choices}
        \CorrectChoice Introducción, ronda de participación y conclusión.
        \choice Introducción, periodo de descanso y conclusión.
        \choice Ronda de participación activa, receso de los ponentes y conclusión general.
        \choice Inicio de la sesión, conversación abierta y conclusión sumaria.
    \end{choices}
    \vspace{0.3cm}

    \question Indica cuáles son los elementos clave que se emplean en el diseño de anuncios bidimensionales:
    \begin{choices}
        \choice Encabezado llamativo, leyenda aclaratoria pequeña y dibujos lineales a mano.
        \choice Dibujo técnico, información estadística, logotipo corporativo y fecha de impresión.
        \CorrectChoice Encabezado, cuerpo del texto, imagen principal y datos precisos del emisor.
        \choice Cuerpo del mensaje, imagen central descriptiva, leyenda de derechos y datos comerciales del emisor.
    \end{choices}
    \vspace{0.3cm}

    \question ¿Cómo están conformados y organizados estructuralmente los folletos dípticos?
    \begin{choices}
        \choice Por 3 partes bien definidas: una portada informativa, el contenido interior y una página final de cierre.
        \choice Por 2 partes sencillas: una portada exterior y una contraportada limpia de texto.
        \choice Por 4 partes complejas: portada, contraportada, bloque de información principal y un resumen ejecutivo.
        \CorrectChoice Por 4 caras bien delimitadas: la portada de presentación, la página 1 interior, la página 2 interior y la contraportada.
    \end{choices}
    \vspace{0.3cm}

    \question ¿Cuál es la diferencia operativa y técnica principal entre los programas de radio tradicionales y los podcasts?
    \begin{choices}
        \choice La radio se escucha exclusivamente en aparatos receptores analógicos, mientras que el podcast corre solo en navegadores de internet de escritorio.
        \CorrectChoice La radio se transmite en vivo en un horario y frecuencia fija, mientras que el podcast puede ser en vivo o grabado y se escucha bajo demanda a elección del usuario.
        \choice La radio fue diseñada únicamente para dar boletines informativos oficiales, mientras que el podcast tiene el fin exclusivo de entretener al oyente.
        \choice La radio es consumida primordialmente por audiencias de edad avanzada, mientras que el formato podcast es de uso exclusivo de audiencias jóvenes.
    \end{choices}
    				    \subsection{Lengua extranjera}

\question Which of the following options represents the correct passive voice form of the active sentence: \textit{"The committee is reviewing the proposal right now"}?
    \begin{choices}
        \choice The proposal is reviewed by the committee right now.
        \choice The proposal was being reviewed by the committee right now.
        \CorrectChoice The proposal is being reviewed by the committee right now.
        \choice The proposal has been reviewed by the committee right now.
    \end{choices}
    \vspace{0.3cm}

    \question Complete the sentence with the correct phrasal verb: \textit{"The manager decided to \underline{\hspace{2cm}} the meeting until next week due to the CEO's unexpected absence."}
    \begin{choices}
        \CorrectChoice put off
        \choice call off
        \choice turn down
        \choice bring up
    \end{choices}
    \vspace{0.3cm}

    \question Choose the correct reported speech (indirect speech) version of the following quote: \textit{"I will call you tomorrow," said Mary.}
    \begin{choices}
        \choice Mary said that she will call me tomorrow.
        \CorrectChoice Mary said that she would call me the next day.
        \choice Mary said that she would have called me tomorrow.
        \choice Mary said that she calls me the following day.
    \end{choices}
    \vspace{0.3cm}

    \question Complete the following third conditional sentence correctly: \textit{"If they had left the house earlier, they \underline{\hspace{2cm}} the flight to London."}
    \begin{choices}
        \choice wouldn't miss
        \choice hadn't missed
        \CorrectChoice wouldn't have missed
        \choice won't have missed
    \end{choices}
    \vspace{0.3cm}

    \question Choose the grammatically correct sentence that correctly applies negative inversion for emphasis:
    \begin{choices}
        \choice Rarely I have seen such a beautiful and majestic sunset.
        \CorrectChoice Rarely have I seen such a beautiful and majestic sunset.
        \choice Have I rarely seen such a beautiful and majestic sunset.
        \choice Rarely saw I such a beautiful and majestic sunset.
    \end{choices}
              
    \newpage

\section{Saberes y pensamiento científico}
\subsection{Matemáticas}

\question En un zoológico hay avestruces y cebras. Si en total se cuentan 35 cabezas y 110 patas, ¿cuántas cebras hay en el recinto?
    \begin{choices}
        \choice 15
        \CorrectChoice 20
        \choice 25
        \choice 30
    \end{choices}

    \question Al lanzar simultáneamente dos dados cúbicos justos de seis caras, ¿cuál es la probabilidad de que la suma de los números que caigan en la cara superior sea un número primo?
    \begin{choices}
        \choice $5/18$
        \choice $1/3$
        \CorrectChoice $5/12$
        \choice $1/2$
    \end{choices}

    \question ¿Cuánto mide la suma de los ángulos interiores de un polígono regular que tiene 20 diagonales en total?
    \begin{choices}
        \choice 720°
        \CorrectChoice 1080° 
        \choice 1440°
        \choice 1800°
    \end{choices}

    \question Al simplificar la expresión algebraica $(2x - 3y)^2 - (2x + 3y)^2$, se obtiene:
    \begin{choices}
        \choice $0$
        \choice $-18y^2$
        \choice $24xy$
        \CorrectChoice $-24xy$
    \end{choices}

    \question Si $3^x = 81$ y $2^{x-y} = 8$, ¿cuál es el valor numérico de la expresión $x^2 - y^2$?
    \begin{choices}
        \choice 7
        \CorrectChoice 15 
        \choice 25
        \choice 31
    \end{choices}

\subsection{Física}

\question Un bloque de 10 kg se empuja sobre una superficie horizontal con una fuerza constante de 50 N. Si la fuerza de fricción cinética entre el bloque y el suelo es de 20 N, ¿cuál es la aceleración del bloque?
    \begin{choices}
        \choice 2 m/s²
        \CorrectChoice 3 m/s²
        \choice 5 m/s²
        \choice 7 m/s²
    \end{choices}

        \newpage

    \question Si un automóvil viaja en línea recta y duplica su velocidad, ¿qué sucede con su energía cinética?
    \begin{choices}
        \choice Se mantiene igual debido a la inercia.
        \choice Se reduce a la mitad.
        \choice Se duplica.
        \CorrectChoice Se cuadruplica.
    \end{choices}

    \question El equivalente mecánico del calor, demostrado experimentalmente por James Joule, estableció formalmente que:
    \begin{choices}
        \choice La temperatura es proporcional a la densidad del material.
        \CorrectChoice El trabajo mecánico y el calor son manifestaciones equivalentes de la transferencia de energía.
        \choice Toda máquina térmica puede alcanzar un 100\% de eficiencia.
        \choice El calor fluye espontáneamente de los cuerpos fríos a los calientes.
    \end{choices}

    \question En un circuito eléctrico simple que obedece la Ley de Ohm, si se cuadruplica el voltaje aplicado a una resistencia constante, la corriente que fluye por el circuito:
    \begin{choices}
        \choice Disminuye a la cuarta parte.
        \choice Se duplica.
        \CorrectChoice Se cuadruplica.
        \choice Permanece inalterada.
    \end{choices}

    \question Una palanca de primer género está en equilibrio rotacional. Si el brazo de palanca de la fuerza de entrada (potencia) es el triple que el brazo de palanca de la carga (resistencia), la ventaja mecánica ideal es:
    \begin{choices}
        \choice 0.33
        \choice 1
        \CorrectChoice 3
        \choice 9
    \end{choices}

    \question Dos cargas eléctricas puntuales, separadas a una distancia $r$, se repelen con una fuerza electrostática $F$. De acuerdo con la Ley de Coulomb, si la distancia entre ellas se triplica (3r), la nueva fuerza de repulsión será:
    \begin{choices}
        \choice $F/3$
        \choice $3F$
        \choice $F/6$
        \CorrectChoice $F/9$
    \end{choices}

    \question El principio de inducción electromagnética, descubierto por Michael Faraday, es el fundamento tecnológico para el funcionamiento de:
    \begin{choices}
        \choice Las baterías químicas alcalinas.
        \CorrectChoice Los generadores eléctricos y los transformadores de corriente alterna.
        \choice Los tubos de vacío y diodos de estado sólido.
        \choice Los paneles solares fotovoltaicos.
    \end{choices}

        \newpage

    \question Al pasar un rayo de luz monocromática desde el aire hacia el interior de un bloque de vidrio transparente con un ángulo de incidencia oblicuo, ¿qué parámetro ondulatorio de la luz cambia, originando la refracción?
    \begin{choices}
        \choice Su frecuencia.
        \choice Su periodo de oscilación.
        \CorrectChoice Su velocidad de propagación.
        \choice Su naturaleza de fotón a onda mecánica.
    \end{choices}

    \question Un isótopo radiactivo tiene una vida media (semivida) de 8 días. Si una muestra médica pura contiene inicialmente 40 gramos del isótopo, ¿qué cantidad permanecerá sin decaer después de 24 días?
    \begin{choices}
        \choice 10 gramos.
        \CorrectChoice 5 gramos. 
        \choice 2.5 gramos.
        \choice 0 gramos.
    \end{choices}

    \question En un circuito eléctrico, tres resistencias idénticas de 12 ohmios se conectan en paralelo. ¿Cuál es la resistencia equivalente total del circuito?
    \begin{choices}
        \choice 36 ohmios.
        \choice 12 ohmios.
        \choice 8 ohmios.
        \CorrectChoice 4 ohmios. 
    \end{choices}

\section{De lo Humano a lo comunitario}
\subsection{Tecnología}

\question ¿Qué son y qué función desempeñan programas como Microsoft PowerPoint y Google Slides?
    \begin{choices}
        \CorrectChoice Son programas que permiten crear presentaciones digitales con texto, imágenes, videos, gráficos y animaciones para exponer información de manera visual y organizada.
        \choice Son navegadores web de alta velocidad desarrollados para la descarga segura de paqueterías de oficina.
        \choice Son sistemas operativos móviles que administran las tareas secundarias del hardware de las tabletas escolares.
        \choice Son lenguajes de programación especializados en el desarrollo estructurado de bases de datos numéricas.
    \end{choices}

    \question ¿Cómo se crea una nueva diapositiva dentro de una presentación de PowerPoint o Google Slides?
    \begin{choices}
        \choice Reiniciando el programa de edición o cerrando el documento de trabajo actual para limpiar el caché.
        \choice Presionando el botón de encendido del monitor o mediante el menú general de configuración del sistema operativo.
        \CorrectChoice Seleccionando la opción ``Nueva diapositiva'' en la barra de herramientas o usando el atajo correspondiente del teclado.
        \choice Insertando un hipervínculo que apunte de forma directa a una página web externa.
    \end{choices}

        \newpage

    \question ¿Qué herramientas visuales y estructurales se pueden utilizar para mejorar sustancialmente una presentación digital?
    \begin{choices}
        \choice Fórmulas algebraicas complejas, macros programadas y hojas de cálculo embebidas de alta densidad.
        \choice Herramientas de edición de código HTML y hojas de estilo en cascada externas.
        \CorrectChoice Temas, diseños, imágenes, transiciones, animaciones, gráficos, colores y diferentes tipos de letra.
        \choice Comandos de consola de administración del sistema y archivos ejecutables del sistema operativo.
    \end{choices}

    \question ¿Cuál es la función específica de las transiciones y de las animaciones en una presentación?
    \begin{choices}
        \choice Las transiciones modifican el tamaño del archivo, mientras que las animaciones exportan la presentación a un formato de video comprimido.
        \CorrectChoice Las transiciones permiten cambiar de una diapositiva a otra con efectos visuales, mientras que las animaciones sirven para dar movimiento a objetos dentro de una misma diapositiva.
        \choice Las transiciones alteran los tipos de letra utilizados, mientras que las animaciones corrigen automáticamente las faltas de ortografía.
        \choice Las transiciones sirven para dar luz de fondo al documento, mientras que las animaciones regulan el volumen del sonido insertado.
    \end{choices}

    \question ¿Por qué resulta fundamental organizar de manera clara la información dentro de una presentación?
    \begin{choices}
        \choice Porque reduce el consumo total de energía eléctrica del proyector o pantalla digital durante la exposición.
        \choice Porque duplica de forma automática la cantidad de diapositivas disponibles sin necesidad de escribir más texto.
        \CorrectChoice Porque facilita la comprensión del tema, mantiene la atención del público y ayuda a comunicar las ideas de manera clara, ordenada y profesional.
        \choice Porque es un requisito técnico obligatorio para poder guardar el archivo en formatos compatibles con teléfonos celulares.
    \end{choices}

\section{Ética, Naturaleza y Sociedades}
    \subsection{Formación Cívica y Ética}

\question En la conformación estructural del Estado mexicano, ¿cuáles son sus tres elementos constitutivos esenciales?
    \begin{choices}
        \choice Constitución, Fuerzas Armadas y Sistema de Partidos.
        \CorrectChoice Territorio, Población y Gobierno (Autoridad o Poder Soberano).
        \choice Poder Ejecutivo, Poder Legislativo y Poder Judicial.
        \choice Nación, Religión y Lengua oficial.
    \end{choices}

    \question La función principal del Poder Legislativo a nivel federal en México, depositado en el Congreso de la Unión, es:
    \begin{choices}
        \choice Administrar el presupuesto público para ejecutar obras de infraestructura.
        \choice Juzgar y emitir sentencias inapelables en tribunales de circuito.
        \CorrectChoice Debatir, crear, modificar y derogar leyes, así como aprobar el presupuesto de egresos.
        \choice Representar a México diplomáticamente ante organismos internacionales.
    \end{choices}

        \newpage

    \question En una sociedad democrática, el pluralismo político y cultural se sustenta éticamente en el principio de:
    \begin{choices}
        \choice La asimilación forzada de las minorías.
        \choice La hegemonía de la ideología con mayor poder económico.
        \CorrectChoice La tolerancia activa y el reconocimiento de que la diversidad enriquece a la nación y debe ser protegida.
        \choice El relativismo moral absoluto, donde ninguna ley es vinculante.
    \end{choices}

    \question La Suprema Corte de Justicia de la Nación (SCJN) desempeña un rol fundamental en el Estado de derecho porque su función primordial es:
    \begin{choices}
        \choice Promulgar nuevas leyes cuando el Congreso no llega a un acuerdo.
        \CorrectChoice Fungir como el máximo tribunal constitucional, garantizando que ninguna ley o acto de autoridad viole la Constitución.
        \choice Investigar los delitos de fuero común cometidos en la vía pública.
        \choice Recaudar los impuestos y sancionar el fraude fiscal directamente.
    \end{choices}

    \question La Comisión Nacional de los Derechos Humanos (CNDH) de México es un organismo público autónomo que tiene la facultad de:
    \begin{choices}
        \choice Emitir sentencias de cárcel contra servidores públicos corruptos.
        \CorrectChoice Investigar presuntas violaciones a derechos humanos y emitir recomendaciones públicas no vinculantes a las autoridades.
        \choice Legislar en materia de equidad de género.
        \choice Modificar la Constitución para incluir nuevos derechos humanos de cuarta generación.
    \end{choices}

    \subsection{Historia}
\question Durante la Ilustración, el filósofo Montesquieu propuso en su obra ``El espíritu de las leyes'' un modelo político diseñado para evitar la tiranía. ¿En qué consistía dicho modelo?
    \begin{choices}
        \choice En la instauración de una república comunitaria sin líderes electos.
        \choice En el fortalecimiento del absolutismo monárquico pero ilustrado.
        \CorrectChoice En la división del poder del Estado en tres ramas autónomas: Ejecutivo, Legislativo y Judicial.
        \choice En la cesión total de la soberanía al pueblo mediante democracia directa diaria.
    \end{choices}

    \question El proceso histórico de la Independencia de las 13 Colonias norteamericanas influyó en la Revolución Francesa primordialmente porque:
    \begin{choices}
        \CorrectChoice Demostró en la práctica que los principios ilustrados de libertad, igualdad y gobierno republicano eran materializables.
        \choice Desvió la atención militar del imperio británico, dejando desprotegido el territorio francés.
        \choice Estableció un bloqueo económico a Europa que provocó la hambruna de París en 1789.
        \choice Obligó al rey Luis XVI a exiliarse en el nuevo continente.
    \end{choices}

        \newpage

    \question ¿Cuál fue el impacto demográfico y geográfico más significativo de la Primera Revolución Industrial en Inglaterra a finales del siglo XVIII?
    \begin{choices}
        \choice La dispersión de la población hacia zonas rurales buscando tierras de cultivo orgánico.
        \CorrectChoice El éxodo rural masivo hacia las ciudades, consolidando la urbanización y el nacimiento de barrios obreros.
        \choice El descenso drástico de la tasa de natalidad debido al trabajo infantil.
        \choice La migración total de la clase burguesa hacia las colonias americanas.
    \end{choices}

    \question En el contexto del México independiente, la Constitución de 1824 estableció como forma de gobierno:
    \begin{choices}
        \choice Una monarquía constitucional moderada liderada por un príncipe europeo.
        \CorrectChoice Una república representativa, popular y federal, dividiendo al país en estados libres y soberanos.
        \choice Una dictadura centralista controlada exclusivamente por el clero secular.
        \choice Una república centralista organizada en departamentos bajo mando militar.
    \end{choices}

    \question La región cultural de Mesoamérica se caracterizó por su diversidad, sin embargo, compartía elementos civilizatorios clave. ¿Cuál de los siguientes no era un rasgo cultural común de Mesoamérica?
    \begin{choices}
        \choice La construcción de basamentos piramidales escalonados.
        \choice El uso del calendario ritual de 260 días y el civil de 365 días.
        \CorrectChoice El empleo extensivo de la metalurgia del hierro para forjar armas de combate.
        \choice El cultivo de la triada agrícola: maíz, frijol y calabaza.
    \end{choices}

\end{questions}
\end{document}
